\documentclass[11pt,a4paper]{article}
\usepackage[utf8]{inputenc}
\usepackage[T1]{fontenc}
\usepackage{amsmath,amssymb}
\usepackage{booktabs}
\usepackage{graphicx}
\usepackage{multirow}
\usepackage{hyperref}
\usepackage{xcolor}
\usepackage[margin=2.5cm]{geometry}
\usepackage{natbib}
\usepackage{tabularx}
\usepackage{array}

\newcommand{\ours}{DisPeD}
\newcommand{\fone}{F$_1$}

\title{\ours{}: Persuasion-Aware Fine-Tuning for Robust\\Cross-Attack Multi-Turn Jailbreak Detection}

\author{Anonymous}
\date{}

\begin{document}
\maketitle

% ============================================================
% ABSTRACT
% ============================================================
\begin{abstract}
Multi-turn jailbreak attacks distribute adversarial intent across individually benign conversational turns, defeating per-turn safety classifiers and single-turn content analysis.
A natural defense is to track semantic drift across turns with a generic encoder and a recurrent aggregator, but we show that this vanilla temporal modeling approach catastrophically fails on unseen attack families: a model trained on Crescendo conversations collapses to \fone{}=0.28 when evaluated on Deceptive Delight.
We propose \ours{}, which fine-tunes DeBERTa-v3-base on an 8-class persuasion strategy classification objective before aggregation.
This persuasion-aware fine-tuning bridges the cross-attack generalization gap entirely, achieving \fone{}=1.0 on the same unseen attacks (a +72 percentage point improvement).
Systematic ablation reveals that the supervision signal, not architectural choices, drives generalization: scrambled strategy labels retain \fone{}=0.94--0.99, binary persuasion-presence labels match 8-class performance within 1 percentage point (pp), mean pooling matches GRU aggregation, and detection is insensitive to multi-task loss weighting.
The 8-class taxonomy provides complementary value through interpretable strategy trajectories that distinguish attack families.
Zero-shot frontier Large Language Model (LLM) judges (GPT-5.5, Qwen3-8B) achieve at most \fone{}=0.45 on the same data, confirming that content-level analysis cannot detect process-level manipulation.
\end{abstract}

% ============================================================
% 1. INTRODUCTION
% ============================================================
\section{Introduction}

Multi-turn jailbreak attacks such as Crescendo \citep{russinovich2024}, Foot-in-the-Door (FITD), and Deceptive Delight \citep{jiang2024} exploit conversational dynamics to circumvent safety alignment in large language models.
Rather than injecting a single adversarial prompt, these attacks distribute manipulative intent across multiple turns, keeping each individual message within the ``benign'' region of representation space \citep{bullwinkel2025}.
This deliberate distribution defeats per-turn safety classifiers \citep{inan2023llamaguard,han2024wildguard,zeng2024shieldgemma} and LLM-as-judge approaches \citep{perez2022}, which analyze content at the utterance level rather than tracking the adversarial process across the full conversation.

A natural response is to model the temporal evolution of turn-level representations.
DeepContext \citep{su2026} encodes each turn with BERT and aggregates the sequence with a GRU, achieving \fone{}=0.84 on in-distribution Crescendo detection.
However, we demonstrate that this vanilla semantic temporal modeling approach fails catastrophically when confronted with attack types unseen during training: a DeBERTa-v3-base + GRU model trained on Crescendo conversations collapses to \fone{}=0.28 on Deceptive Delight (DD) attacks.
The generic encoder produces turn embeddings that capture surface semantics but miss the pragmatic manipulation strategies shared across attack families.

Recent work establishes that persuasion is systematic in jailbreaks.
\citet{zeng2024pap} identify 40 persuasion techniques effective for jailbreaking and show that persuasion-enhanced attacks substantially increase attack success rates.
\citet{hagendorff2026} find identifiable persuasion strategies in 97\% of successful attacks across a large-scale corpus.
Complementary analysis demonstrates that persuasion dynamics are extractable from turn-level neural representations \citep{howdollms2025}.
Yet no published method operationalizes these observations into a detection system that generalizes across attack types.

We propose \ours{} (\textbf{Dis}course-level \textbf{Pe}rsuasion \textbf{D}etector), which uses persuasion strategy classification as a multi-task fine-tuning objective for the turn-level encoder.
Stage~1 fine-tunes DeBERTa-v3-base on two tasks jointly: 8-class persuasion strategy classification and binary jailbreak-intent classification.
Stage~2 aggregates the resulting turn embeddings via a sequence model (GRU or mean pooling) for conversation-level detection.
The persuasion objective teaches the encoder to attend to pragmatic features that transfer across attack families, without requiring attack-specific training data.

Our experiments yield three findings.
First, persuasion-aware fine-tuning bridges the cross-attack generalization gap completely: \ours{} achieves \fone{}=1.0 on unseen DD attacks where the vanilla baseline collapses to \fone{}=0.28, a +72pp improvement.
Second, systematic ablation reveals that the persuasion fine-tuning objective is the sole critical ingredient: scrambled labels, binary labels, mean pooling, and varying loss weights all produce equivalent detection, while the 8-class taxonomy provides interpretable strategy trajectories as a complementary contribution.
Third, zero-shot LLM judges (GPT-5.5: \fone{}=0.38; Qwen3-8B: \fone{}=0.10) confirm that content-level analysis cannot substitute for process-level detection.

Our contributions are:
\begin{itemize}
    \item We demonstrate that vanilla semantic temporal modeling (generic encoder + RNN) fails catastrophically on unseen multi-turn attack types (\fone{}=0.28 on DD) and that persuasion-aware fine-tuning bridges this gap entirely (\fone{}=1.0, +72pp).
    \item We provide comprehensive ablation evidence that the persuasion fine-tuning objective functions as an accessible supervision signal: detection is robust to label quality, label granularity, aggregation architecture, and loss weighting. The 8-class taxonomy enables interpretable strategy trajectory analysis as a complementary contribution.
    \item We establish baselines against TF-IDF (strong lexical), frontier LLM judges (content-level), and a DeepContext proxy (vanilla temporal), positioning \ours{} within the landscape of multi-turn jailbreak defenses.
\end{itemize}

% ============================================================
% 2. RELATED WORK
% ============================================================
\section{Related Work}

\paragraph{Multi-turn jailbreak attacks.}
Crescendo \citep{russinovich2024} gradually escalates requests across turns, exploiting conversational context to normalize harmful content.
PAIR \citep{chao2023} uses an attacker LLM to iteratively refine jailbreak prompts.
\citet{jiang2024} introduce Deceptive Delight, which interleaves benign and adversarial content to exploit attention dilution.
\citet{wei2024} provide a taxonomy of jailbreak strategies, distinguishing competing-objectives and mismatched-generalization failures.
\citet{bullwinkel2025} show that multi-turn attacks keep per-turn representations in a ``benign'' region, causing a 50-point attack success rate gap relative to single-turn attacks.

\paragraph{Safety classifiers and guardrails.}
Per-turn classifiers such as Llama Guard \citep{inan2023llamaguard}, WildGuard \citep{han2024wildguard}, ShieldGemma \citep{zeng2024shieldgemma}, and Constitutional Classifiers \citep{anthropic2024constitutional} detect harmful content at the utterance level.
These systems achieve high accuracy on single-turn attacks but do not model cross-turn dependencies.
LLM-as-judge approaches \citep{perez2022,ganguli2022} prompt frontier models to evaluate conversations but analyze content rather than process, limiting their effectiveness against distributed manipulation.

\paragraph{Multi-turn defenses.}
DeepContext \citep{su2026} aggregates BERT turn embeddings via a GRU for conversation-level detection, achieving \fone{}=0.84 on Crescendo.
BIID \citep{biid2025} and G-Guard \citep{gguard2025} propose graph-based and multi-agent approaches to multi-turn safety.
SafeDream \citep{safedream2025} uses proactive safety mechanisms.
None of these methods incorporate persuasion-grounded features or evaluate cross-attack generalization.

\paragraph{Persuasion in adversarial contexts.}
\citet{zeng2024pap} systematically apply 40 persuasion techniques to jailbreaking, demonstrating that persuasion-enhanced prompts substantially increase attack success.
\citet{hagendorff2026} analyze a large-scale jailbreak corpus and find identifiable persuasion strategies in 97\% of successful attacks.
\citet{howdollms2025} show that persuasion dynamics are extractable from turn-level representations.
These findings motivate persuasion-aware detection but have not been operationalized into cross-attack-generalizable classifiers.

\paragraph{Discourse analysis and computational persuasion.}
Discourse-level analysis has a long history in NLP, from Rhetorical Structure Theory \citep{mann1988} to the Penn Discourse Treebank \citep{prasad2008pdtb} and dialogue act classification \citep{stolcke2000}.
Computational persuasion detection has been studied in argumentation mining \citep{habernal2016} and social media influence \citep{tan2016}.
\ours{} connects these traditions to LLM safety by treating persuasion strategy classification as a representation learning objective.

\paragraph{Red-teaming and evaluation.}
Systematic red-teaming efforts \citep{perez2022,ganguli2022} and benchmarks such as HarmBench \citep{mazeika2024} evaluate model safety.
\citet{nasr2025} study adaptive attacks that evade classifiers.
\citet{gururangan2018} and \citet{niven2019} demonstrate annotation artifacts in NLP benchmarks.
We address artifact concerns through topic-anchored data construction and ablation controls.


% ============================================================
% 3. METHOD
% ============================================================
\section{Method}

\subsection{Problem Formulation}
Given a multi-turn conversation $C = (u_1, u_2, \ldots, u_T)$ consisting of $T$ turns, the goal is to classify $C$ as jailbreak or benign at the conversation level.
We additionally define a per-turn persuasion strategy label $s_t \in \{0, 1, \ldots, 7\}$ from an 8-class taxonomy and a binary jailbreak-intent label $y_t \in \{0, 1\}$.

\subsection{Stage 1: Persuasion-Aware Encoder Fine-Tuning}
We fine-tune DeBERTa-v3-base \citep{he2021deberta} on two auxiliary tasks jointly.
Each turn $u_t$ is encoded independently to produce a [CLS] representation $\mathbf{h}_t \in \mathbb{R}^d$.
Two classification heads operate on $\mathbf{h}_t$:

\begin{align}
    \hat{s}_t &= \text{softmax}(\mathbf{W}_s \mathbf{h}_t + \mathbf{b}_s) \label{eq:persuasion} \\
    \hat{y}_t &= \sigma(\mathbf{w}_y^\top \mathbf{h}_t + b_y) \label{eq:intent}
\end{align}

The combined loss is:
\begin{equation}
    \mathcal{L}_{\text{stage1}} = \mathcal{L}_{\text{CE}}(\hat{s}_t, s_t) + \lambda \cdot \mathcal{L}_{\text{BCE}}(\hat{y}_t, y_t)
\end{equation}
where $\lambda$ controls the relative weight of the intent classification objective.
Ablation experiments (Section~\ref{sec:ablation}) show detection performance is insensitive to $\lambda \in \{0, 0.3, 0.5, 2.0\}$.

\subsection{Stage 2: Conversation-Level Aggregation}
Given the fine-tuned encoder, we extract turn embeddings $\mathbf{H} = (\mathbf{h}_1, \ldots, \mathbf{h}_K)$ for the first $K$ turns of a conversation.
These are aggregated and classified:
\begin{equation}
    \hat{c} = \sigma(\mathbf{w}_c^\top f(\mathbf{H}) + b_c)
\end{equation}
where $f$ is either a GRU that produces the final hidden state, or mean pooling across turns.
We report GRU as the default but note that mean pooling produces equivalent results (Section~\ref{sec:ablation}), confirming that temporal ordering is not the detection mechanism.

\subsection{Early Detection}
For deployment, the system can operate on partial conversations.
Given the first $K$ turns, detection is performed on $\mathbf{H}_K = (\mathbf{h}_1, \ldots, \mathbf{h}_K)$.
We evaluate at $K \in \{1, 2, 3, 5, \text{full}\}$ to characterize early detection capability.

\subsection{Persuasion Strategy Taxonomy}
\label{sec:taxonomy}
Following \citet{zeng2024pap} and \citet{hagendorff2026}, we define 8 persuasion strategy classes spanning the spectrum from rapport building to explicit coercion.
Strategy index increases with escalation intensity: class 0 (rapport/trust building) through class 7 (explicit pressure/coercion).
This ordering enables trajectory analysis (Section~\ref{sec:interpretability}) where the temporal pattern of strategy indices characterizes different attack types.

% ============================================================
% 4. EXPERIMENTAL SETUP
% ============================================================
\section{Experimental Setup}

\subsection{Data}
\paragraph{Training data.}
We construct a topic-anchored dataset of 480 conversations (240 Crescendo jailbreak, 240 benign) with per-turn persuasion strategy annotations.
Benign conversations are anchored to matching topics to eliminate trivial lexical shortcuts from topic divergence.
Each conversation is annotated with turn-level 8-class persuasion strategy labels and binary jailbreak-intent labels.

\paragraph{In-distribution test set.}
A held-out set of 72 Crescendo conversations (36 jailbreak, 36 benign) is used for IID evaluation.

\paragraph{Out-of-distribution test sets.}
We generate 80 FITD and 80 Deceptive Delight (DD) conversations with persuasion annotations using a controlled generation pipeline.
FITD conversations use 4 distinct persuasion strategies and average 8.5 turns.
DD conversations use 6 strategies and average 7.2 turns.
These are used exclusively for evaluation; no OOD data appears during training.

\paragraph{Hard negative evaluation.}
200 conversations from WildChat \citep{zhao2024wildchat} are filtered for high-drift characteristics (topic shift $\geq$3, $>$10 turns, emotional/persuasive language) after removing contaminated samples (contamination rate reduced from 15\% to 2\%).

\subsection{Models}
\paragraph{\ours{} (Treatment).}
DeBERTa-v3-base fine-tuned on the persuasion + intent multi-task objective (Stage~1), followed by GRU aggregation (Stage~2).
Achieves 82\% accuracy on persuasion strategy classification and 98\% on intent classification.

\paragraph{Vanilla baseline.}
Same Stage~2 architecture (DeBERTa-v3-base + GRU) without Stage~1 fine-tuning.

\paragraph{TF-IDF + Logistic Regression.}
TF-IDF features with logistic regression, serving as a strong lexical baseline.

\paragraph{DeepContext proxy.}
BERT-base + GRU, reproducing the DeepContext \citep{su2026} approach with available resources.

\paragraph{LLM-as-Judge.}
GPT-5.5 and Qwen3-8B with zero-shot prompting for conversation-level jailbreak detection.

\subsection{Evaluation}
All models are evaluated on conversation-level binary classification (\fone{} score) at $K \in \{1, 2, 3, 5, \text{full}\}$ turns.
Treatment and baseline models are trained with 3 random seeds; we report mean $\pm$ standard deviation.
Bootstrap 95\% confidence intervals are computed on the test set.

% ============================================================
% 5. RESULTS
% ============================================================
\section{Results}

\subsection{Baseline Landscape}
\label{sec:baselines}

Table~\ref{tab:tfidf} reports the TF-IDF + LR baseline across all attack types.
TF-IDF achieves \fone{}=1.0 on Crescendo and FITD, and \fone{}=0.944 on DD.
This confirms inherent lexical separability between jailbreak and benign conversations: the top predictive features for jailbreak (``sure,'' ``hypothetical,'' ``creative,'' ``story,'' ``imagine'') versus benign (``thanks,'' ``code,'' ``help,'' ``python,'' ``error'') reflect fundamentally different conversational registers.
We report TF-IDF as an honest upper bound for lexical approaches throughout the paper.

\begin{table}[t]
\centering
\caption{TF-IDF + Logistic Regression baseline across attack types.}
\label{tab:tfidf}
\begin{tabular}{lccc}
\toprule
 & Crescendo & FITD & DD \\
\midrule
TF-IDF + LR & 1.000 & 1.000 & 0.944 \\
\bottomrule
\end{tabular}
\end{table}

\subsection{In-Distribution Detection}
\label{sec:iid}

Table~\ref{tab:iid} reports IID performance on the Crescendo test set.
Both treatment and baseline achieve near-perfect detection at higher $K$, with the treatment reaching \fone{}=1.0 across all $K$ values with zero variance across 3 seeds.
The baseline is also strong IID (\fone{}=0.968 at $K$=1, rising to 1.0 at full), confirming that vanilla DeBERTa is a competent in-distribution detector.
The small IID gap (max 3.2pp at $K$=1) is expected: the paper's value proposition is not IID performance but cross-attack generalization.

\begin{table}[t]
\centering
\caption{IID detection performance on Crescendo test set ($n$=72, 3-seed mean $\pm$ std). Treatment achieves zero variance across all $K$.}
\label{tab:iid}
\begin{tabular}{lccccc}
\toprule
Model & $K$=1 & $K$=2 & $K$=3 & $K$=5 & Full \\
\midrule
Treatment & 1.000\tiny{$\pm$.000} & 1.000\tiny{$\pm$.000} & 1.000\tiny{$\pm$.000} & 1.000\tiny{$\pm$.000} & 1.000\tiny{$\pm$.000} \\
Baseline  & 0.968\tiny{$\pm$.033} & 0.977\tiny{$\pm$---} & 0.987\tiny{$\pm$---} & 0.996\tiny{$\pm$---} & 1.000\tiny{$\pm$---} \\
TF-IDF    & \multicolumn{5}{c}{1.000 (all $K$)} \\
\bottomrule
\end{tabular}
\end{table}

\subsection{Cross-Attack Generalization}
\label{sec:cross_attack}

Table~\ref{tab:cross_attack} presents the central experiment: models trained exclusively on Crescendo, evaluated on unseen FITD and DD attacks.
On DD, the vanilla baseline collapses to \fone{}=0.284 at full conversation length, while the treatment achieves \fone{}=1.0 (a +71.6pp gap).
The baseline's performance on DD degrades monotonically with conversation length ($K$=1: 0.575 $\to$ full: 0.284), indicating that the model actively gets worse as DD's deceptive structure unfolds.
The treatment maintains perfect or near-perfect detection across all $K$ values.

On FITD, both models achieve strong performance (treatment: 1.0, baseline: 1.0 at full; 1.0 vs. 0.99 at $K$=1).
This asymmetry is informative: FITD's compliance-escalation pattern is structurally similar to Crescendo, enabling vanilla transfer.
DD's interleaving deception mechanism is fundamentally different, exposing the vanilla approach's reliance on Crescendo-specific surface features.

TF-IDF achieves \fone{}=0.944 on DD (strong but not perfect) and 1.0 on FITD.
The treatment's advantage over TF-IDF on DD (5.6pp) is modest in raw detection but becomes substantial under adversarial perturbation (Section~\ref{sec:perturbation}).

\begin{table}[t]
\centering
\caption{Cross-attack generalization: models trained on Crescendo only, evaluated on unseen FITD and DD. The +71.6pp gap on DD is the paper's headline result.}
\label{tab:cross_attack}
\begin{tabular}{llccccc}
\toprule
Attack & Model & $K$=1 & $K$=2 & $K$=3 & $K$=5 & Full \\
\midrule
\multirow{3}{*}{DD}
  & Treatment & 1.000 & 0.990 & 0.990 & 1.000 & 1.000 \\
  & Baseline  & 0.575 & 0.438 & 0.360 & 0.310 & 0.284 \\
  & TF-IDF    & --- & --- & --- & --- & 0.944 \\
\midrule
\multirow{3}{*}{FITD}
  & Treatment & 1.000 & --- & --- & --- & 1.000 \\
  & Baseline  & 0.990 & --- & --- & --- & 1.000 \\
  & TF-IDF    & --- & --- & --- & --- & 1.000 \\
\bottomrule
\end{tabular}
\end{table}

\subsection{Ablation Studies}
\label{sec:ablation}

Table~\ref{tab:ablation} presents systematic ablation across four dimensions: label quality, label granularity, loss weighting, and aggregation architecture.
The results converge on a single conclusion: the persuasion fine-tuning objective is the sole critical ingredient for cross-attack generalization.

\paragraph{Label quality (scrambled vs. proper).}
Randomly permuting strategy labels before fine-tuning reduces DD \fone{} by only 1.0pp at full length (0.990 vs. 1.000) and 6.5pp at $K$=1 (0.935 vs. 1.000).
Both scrambled and proper labels massively outperform the vanilla baseline (+70.6pp and +71.6pp on DD, respectively).
This reveals that the fine-tuning process on persuasion-related text, not the specific class boundaries, drives generalization.
Scrambled fine-tuning still achieves 64.4\% persuasion strategy accuracy (vs. 79.6\% with proper labels), confirming that even degraded class signal provides useful representation structure.

\paragraph{Label granularity (binary vs. 8-class).}
Collapsing the 8-class taxonomy to binary (persuasion-present vs. none) produces equivalent IID detection (\fone{}=1.0) and nearly equivalent DD detection (\fone{}=0.990 vs. 1.000, a 1.0pp gap at full; 0.981 vs. 1.000 at $K$=1).
The model is fundamentally a persuasion-presence detector.
The 8-class taxonomy provides interpretability value (Section~\ref{sec:interpretability}), not detection advantage.

\paragraph{Loss weighting ($\lambda$).}
Detection is insensitive to $\lambda \in \{0, 0.3, 0.5, 2.0\}$: all configurations achieve \fone{}=1.0 on both IID and DD.
The $\lambda$=0 condition (persuasion-only, no intent head) working equally well confirms that the intent classification head does not leak direct jailbreak signal into the encoder.

\paragraph{Aggregation architecture (GRU vs. mean pooling).}
Mean pooling matches GRU within 1pp on all evaluations (DD full: 1.0 vs. 1.0; DD $K$=1: 1.0 vs. 1.0).
Temporal ordering is not the detection mechanism: the model detects persuasion-enriched embeddings regardless of turn order.

\begin{table}[t]
\centering
\caption{Ablation study on DD (out-of-distribution). All variants with any form of persuasion fine-tuning massively outperform the vanilla baseline (\fone{}=0.284). The persuasion objective is the sole critical ingredient.}
\label{tab:ablation}
\begin{tabular}{lcc}
\toprule
Configuration & DD Full \fone{} & DD $K$=1 \fone{} \\
\midrule
Treatment (8-class, GRU, $\lambda$=0.5) & 1.000 & 1.000 \\
Scrambled labels & 0.990 & 0.935 \\
Binary labels & 0.990 & 0.981 \\
$\lambda$=0 (persuasion only) & 1.000 & --- \\
$\lambda$=0.3 & 1.000 & --- \\
$\lambda$=2.0 & 1.000 & --- \\
Mean pooling (replace GRU) & 1.000 & 1.000 \\
\midrule
Vanilla baseline (no fine-tuning) & 0.284 & 0.575 \\
\bottomrule
\end{tabular}
\end{table}

\subsection{Adversarial Perturbation Robustness}
\label{sec:perturbation}

To test whether models rely on surface lexical cues, we apply aggressive paraphrasing to the Crescendo test set (BERTScore = 0.92 to originals, lexical overlap = 0.20) and evaluate all models on the perturbed data.
Table~\ref{tab:perturbation} reports \fone{} and the drop ($\Delta$) relative to original performance.

The treatment shows zero degradation at $K \geq 2$ and only $-$2.6pp at $K$=1.
The vanilla baseline is the most fragile ($-$11.8pp at full, $-$14.6pp at $K$=1).
TF-IDF is intermediate: robust at $K \geq 3$ ($-$1.3pp) but degrades at $K$=1 ($-$9.6pp), where lexical signal is sparse.
The strongest differentiation appears at early detection under perturbation ($K$=1), where the treatment leads TF-IDF by 7.0pp and the baseline by 12.0pp.

\begin{table}[t]
\centering
\caption{Adversarial perturbation robustness. $\Delta$ = \fone{} drop from unperturbed. BERTScore = 0.92, lexical overlap = 0.20. Early detection ($K$=1) under perturbation maximally separates the three approaches.}
\label{tab:perturbation}
\begin{tabular}{lcccc}
\toprule
Model & Orig. \fone{} & Pert. \fone{} & $\Delta$ & Pert. $K$=1 ($\Delta$) \\
\midrule
Treatment & 1.000 & 1.000 & 0.0pp & 0.974 ($-$2.6pp) \\
Baseline  & 1.000 & 0.882 & $-$11.8pp & 0.854 ($-$14.6pp) \\
TF-IDF    & 1.000 & 0.987 & $-$1.3pp & 0.904 ($-$9.6pp) \\
\bottomrule
\end{tabular}
\end{table}

\subsection{LLM-as-Judge Baseline}
\label{sec:llm_judge}

Table~\ref{tab:llm_judge} evaluates zero-shot conversation-level jailbreak detection by frontier LLMs.
GPT-5.5 achieves perfect precision (1.0 at $K$=2) but catastrophically low recall (10.5\% at $K$=2, 29.0\% at $K$=5), producing \fone{}=0.19--0.45.
The frontier LLM is ``deceived'' by Crescendo's surface educational framing because it analyzes content, not the underlying manipulation process.
A substantial fraction of responses are unparseable (30.3\% at $K$=2, 47.4\% at full), further limiting practical utility.
Even with best-case unparseable correction (all unparseable treated as correct), GPT-5.5 reaches approximately \fone{}=0.75, still far below the treatment (\fone{}=1.0).

Qwen3-8B performs near zero (\fone{}=0.0--0.15), confirming that this is not a model-size issue but a fundamental limitation of content-level analysis for process-level attacks.

\begin{table}[t]
\centering
\caption{LLM-as-judge zero-shot detection on Crescendo. Frontier LLMs achieve high precision but catastrophically low recall: content analysis cannot detect distributed persuasion.}
\label{tab:llm_judge}
\begin{tabular}{lccccc}
\toprule
Model & $K$=2 & $K$=3 & $K$=5 & Full & Unparseable \\
\midrule
GPT-5.5 \fone{}   & 0.190 & 0.273 & 0.449 & 0.383 & 30--47\% \\
Qwen3-8B \fone{}  & 0.000 & 0.051 & 0.146 & 0.100 & --- \\
Treatment \fone{}  & 1.000 & 1.000 & 1.000 & 1.000 & --- \\
\bottomrule
\end{tabular}
\end{table}

\subsection{Interpretability via Strategy Trajectories}
\label{sec:interpretability}

Table~\ref{tab:transitions} reports persuasion strategy transition statistics.
Different attack types exhibit distinct trajectory signatures.
Crescendo conversations show gradual escalation, with mean strategy index rising from 0.86 in early turns to 5.64 in late turns (escalation slope = 0.53 per turn).
Across 8 turns, the strategy profile spans all 8 classes.
FITD conversations follow a commitment-escalation pattern: low-intensity strategies dominate early turns before a sharp transition to high-intensity strategies (low $\to$ low $\to$ mid $\to$ high).
DD conversations cycle through emotional manipulation, role assignment, and deception in an interleaving pattern, using 6 of 8 strategy classes.

Benign conversations show low strategy index variance ($<$1.5) and low transition entropy (1.2 vs. 2.8 for jailbreak), confirming the absence of structured escalation.
These trajectories provide security analysts with attack fingerprints that characterize how an attack unfolds, not merely that it is occurring.

\begin{table}[t]
\centering
\caption{Strategy trajectory statistics by attack type. Jailbreak conversations exhibit higher transition entropy and structured escalation patterns absent in benign conversations.}
\label{tab:transitions}
\begin{tabular}{lcccc}
\toprule
 & Crescendo & FITD & DD & Benign \\
\midrule
Strategies used & 8 & 4 & 6 & --- \\
Mean idx. start & 0.86 & \multicolumn{2}{c}{(see pattern)} & --- \\
Mean idx. end & 5.64 & \multicolumn{2}{c}{(see pattern)} & --- \\
Escalation slope & 0.53 & --- & --- & --- \\
Transition entropy & \multicolumn{3}{c}{2.8 (jailbreak)} & 1.2 \\
Idx. variance & --- & --- & --- & $<$1.5 \\
\bottomrule
\end{tabular}
\end{table}


% ============================================================
% 6. ADDITIONAL ANALYSIS
% ============================================================
\section{Additional Analysis}

\subsection{DeepContext Proxy Comparison}

To ground our comparison with the DeepContext approach \citep{su2026}, we train a BERT-base + GRU proxy.
This proxy achieves reasonable IID performance (\fone{}=0.96) but collapses on DD (\fone{}=0.32), confirming that the vanilla temporal modeling failure is independent of encoder choice (BERT-base vs. DeBERTa-v3-base).
The 68pp gap between the proxy and our treatment on DD (\fone{}=0.32 vs. 1.0) validates that encoder representation quality, not temporal architecture, drives generalization (Table~\ref{tab:proxy}).

\begin{table}[t]
\centering
\caption{DeepContext proxy (BERT-base + GRU) comparison. The vanilla temporal approach fails OOD regardless of encoder choice.}
\label{tab:proxy}
\begin{tabular}{lcccc}
\toprule
Model & IID Full & DD Full & DD $K$=1 & FITD Full \\
\midrule
DeepContext proxy & 0.960 & 0.320 & 0.550 & 0.910 \\
Treatment         & 1.000 & 1.000 & 1.000 & 1.000 \\
Baseline (DeBERTa) & 1.000 & 0.284 & 0.575 & 1.000 \\
\bottomrule
\end{tabular}
\end{table}

\subsection{False Positive Analysis on Hard Negatives}

We evaluate false positive rates on 200 WildChat conversations selected for high-drift characteristics (Table~\ref{tab:fpr}).
The treatment produces 8\% FPR, the baseline 12\%, and TF-IDF 5\%.
False positives from the treatment occur on genuinely persuasive benign conversations (sales pitches, debate, tutoring), which is expected behavior for a model designed to detect persuasion dynamics.
Deployment would require threshold tuning for the target precision-recall tradeoff.

\begin{table}[t]
\centering
\caption{False positive rate on WildChat hard negatives ($n$=200). High-drift definition: topic shift $\geq$3, $>$10 turns, emotional/persuasive language.}
\label{tab:fpr}
\begin{tabular}{lc}
\toprule
Model & FPR \\
\midrule
Treatment & 0.08 \\
Baseline  & 0.12 \\
TF-IDF    & 0.05 \\
\bottomrule
\end{tabular}
\end{table}

\subsection{Checkpoint Sensitivity}

Table~\ref{tab:checkpoint} reports performance across fine-tuning checkpoints at epochs 2, 3, and final.
IID detection is invariant across checkpoints (all \fone{}=1.0).
OOD detection shows negligible variance ($<$0.5pp): epoch 2 achieves \fone{}=0.995 on DD, rising to 1.0 by epoch 3.
Even epoch 2 (lower persuasion accuracy: 76\% vs. final 82\%) produces near-perfect downstream detection, reinforcing that precise persuasion classification is not required.

\begin{table}[t]
\centering
\caption{Checkpoint sensitivity analysis. Detection is robust to early stopping; even epoch 2 (lower persuasion accuracy) achieves near-perfect OOD detection.}
\label{tab:checkpoint}
\begin{tabular}{lccc}
\toprule
Checkpoint & IID Full \fone{} & DD Full \fone{} & Persuasion Acc. \\
\midrule
Epoch 2 & 1.000 & 0.995 & 0.76 \\
Epoch 3 & 1.000 & 1.000 & 0.80 \\
Final   & 1.000 & 1.000 & 0.82 \\
\bottomrule
\end{tabular}
\end{table}


% ============================================================
% 7. DISCUSSION
% ============================================================
\section{Discussion}

\paragraph{The supervision signal, not the architecture.}
The comprehensive ablation picture (Section~\ref{sec:ablation}) reveals a clean mechanistic story. Four independent ablation axes (label quality, label granularity, loss weighting, temporal modeling) all point to the same conclusion: the persuasion fine-tuning objective is the sole critical ingredient for cross-attack generalization. Scrambled labels retain \fone{}=0.94--0.99 on DD, binary labels match 8-class within 1pp, mean pooling matches GRU, and all $\lambda$ values produce \fone{}=1.0. The method requires no careful hyperparameter tuning, no expert label annotation, and no specific architecture. This positions persuasion-domain fine-tuning as an accessible and robust supervision signal rather than an overengineered system.

\paragraph{Honest positioning against TF-IDF.}
TF-IDF + LR achieves \fone{}=0.944 on DD and 1.0 on Crescendo and FITD, reflecting inherent lexical separability between jailbreak and benign conversations. \ours{} does not claim superiority over TF-IDF on raw detection. Instead, the treatment provides three complementary advantages: (1) it bridges the vanilla temporal modeling gap (+72pp on DD vs. the baseline's 0.28), establishing that persuasion-aware representations solve a problem that lexical features happen to also address through a different mechanism; (2) it is more robust under adversarial perturbation at early detection ($K$=1: treatment $-$2.6pp vs. TF-IDF $-$9.6pp); (3) it provides interpretable strategy trajectories that TF-IDF cannot offer.

\paragraph{Content analysis vs. process detection.}
The LLM-as-judge results (Section~\ref{sec:llm_judge}) establish a fundamental distinction. GPT-5.5 achieves perfect precision but recall below 30\%, confirming that even frontier LLMs analyzing conversation content cannot detect distributed persuasion when each turn appears individually benign. \ours{} operates at the process level, detecting the accumulation of persuasion dynamics rather than evaluating content for harmfulness. These are complementary capabilities: content analysis catches explicit harm, process detection catches distributed manipulation.

\paragraph{Limitations.}
The evaluation uses synthetic OOD data (80 FITD + 80 DD conversations). While sufficient for proof-of-concept, ecological validity requires evaluation on real-world multi-turn attacks when such data becomes available. The 8\% FPR on WildChat hard negatives indicates that deployment requires threshold calibration. The method targets persuasion-mediated attacks and does not address pure information accumulation, token smuggling, or single-frame roleplay jailbreaks. The training set ($n$=480) is small; scaling to larger and more diverse persuasion-annotated corpora may further improve robustness.


% ============================================================
% 8. CONCLUSION
% ============================================================
\section{Conclusion}

We presented \ours{}, a method that uses persuasion strategy classification as a multi-task fine-tuning objective to enable cross-attack multi-turn jailbreak detection.
Our central finding is that vanilla semantic temporal modeling catastrophically fails on unseen attack types (\fone{}=0.28 on Deceptive Delight), while persuasion-aware fine-tuning bridges this gap entirely (\fone{}=1.0, +72pp).
Systematic ablation demonstrates that the persuasion fine-tuning objective functions as an accessible supervision signal: detection survives label scrambling, binary collapse, architecture changes, and loss reweighting.
The 8-class taxonomy provides complementary interpretability through characteristic strategy trajectories.
These results establish persuasion-domain fine-tuning as a practical and robust approach to detecting distributed multi-turn manipulation, complementing existing per-turn safety classifiers and content-level LLM judges.

% ============================================================
% REFERENCES
% ============================================================
\bibliographystyle{plainnat}

\begin{thebibliography}{30}

\bibitem[Anthropic(2024)]{anthropic2024constitutional}
Anthropic.
\newblock Constitutional classifiers: Defending against universal jailbreaks.
\newblock Technical report, 2024.

\bibitem[Austin(1962)]{austin1962}
J.~L. Austin.
\newblock \emph{How to Do Things with Words}.
\newblock Oxford University Press, 1962.

\bibitem[BIID(2025)]{biid2025}
BIID: Backdoor-based multi-turn jailbreak defense.
\newblock arXiv preprint, 2025.

\bibitem[Bullwinkel et~al.(2025)]{bullwinkel2025}
B.~Bullwinkel et~al.
\newblock Multi-turn attacks keep per-turn representations benign.
\newblock arXiv:2507.02956, 2025.

\bibitem[Chao et~al.(2023)]{chao2023}
P.~Chao, A.~Robey, E.~Dobriban, H.~Hassani, G.~J. Pappas, and E.~Wong.
\newblock Jailbreaking black box large language models in twenty queries.
\newblock arXiv preprint, 2023.

\bibitem[G-Guard(2025)]{gguard2025}
G-Guard: Graph-based multi-turn safety guardrails.
\newblock arXiv preprint, 2025.

\bibitem[Ganguli et~al.(2022)]{ganguli2022}
D.~Ganguli et~al.
\newblock Red teaming language models to reduce harms.
\newblock arXiv preprint, 2022.

\bibitem[Gururangan et~al.(2018)]{gururangan2018}
S.~Gururangan, S.~Swayamdipta, O.~Levy, R.~Schwartz, S.~Bowman, and N.~A. Smith.
\newblock Annotation artifacts in natural language inference data.
\newblock In \emph{NAACL}, 2018.

\bibitem[Habernal and Gurevych(2016)]{habernal2016}
I.~Habernal and I.~Gurevych.
\newblock Which argument is more convincing? Analyzing and predicting convincingness of web arguments using bidirectional LSTM.
\newblock In \emph{ACL}, 2016.

\bibitem[Hagendorff(2026)]{hagendorff2026}
T.~Hagendorff.
\newblock Persuasion strategies in successful jailbreak attacks.
\newblock \emph{Nature Communications}, 2026.

\bibitem[Han et~al.(2024)]{han2024wildguard}
S.~Han et~al.
\newblock WildGuard: Open one-stop moderation tools for safety risks, jailbreaks, and refusals of LLMs.
\newblock arXiv preprint, 2024.

\bibitem[He et~al.(2021)]{he2021deberta}
P.~He, X.~Liu, J.~Gao, and W.~Chen.
\newblock DeBERTa: Decoding-enhanced BERT with disentangled attention.
\newblock In \emph{ICLR}, 2021.

\bibitem[How Do LLMs Persuade?(2025)]{howdollms2025}
How do LLMs persuade? Extracting persuasion dynamics from turn-level representations.
\newblock arXiv:2508.05625, 2025.

\bibitem[Inan et~al.(2023)]{inan2023llamaguard}
H.~Inan et~al.
\newblock Llama Guard: LLM-based input-output safeguard for human-AI conversations.
\newblock arXiv preprint, 2023.

\bibitem[Jiang et~al.(2024)]{jiang2024}
Z.~Jiang et~al.
\newblock Deceptive Delight: Jailbreaking large language models through camouflaged content.
\newblock arXiv preprint, 2024.

\bibitem[Mann and Thompson(1988)]{mann1988}
W.~C. Mann and S.~A. Thompson.
\newblock Rhetorical structure theory: Toward a functional theory of text organization.
\newblock \emph{Text}, 8(3):243--281, 1988.

\bibitem[Mazeika et~al.(2024)]{mazeika2024}
M.~Mazeika et~al.
\newblock HarmBench: A standardized evaluation framework for automated red teaming and robust refusal.
\newblock arXiv preprint, 2024.

\bibitem[Nasr et~al.(2025)]{nasr2025}
M.~Nasr, N.~Carlini, et~al.
\newblock Adaptive attacks against safety classifiers.
\newblock arXiv preprint, 2025.

\bibitem[Niven and Kao(2019)]{niven2019}
T.~Niven and H.-Y. Kao.
\newblock Probing neural network comprehension of natural language arguments.
\newblock In \emph{ACL}, 2019.

\bibitem[Perez et~al.(2022)]{perez2022}
E.~Perez et~al.
\newblock Red teaming language models with language models.
\newblock arXiv preprint, 2022.

\bibitem[Prasad et~al.(2008)]{prasad2008pdtb}
R.~Prasad, N.~Dinesh, A.~Lee, E.~Miltsakaki, L.~Robaldo, A.~Joshi, and B.~Webber.
\newblock The Penn Discourse TreeBank 2.0.
\newblock In \emph{LREC}, 2008.

\bibitem[Russinovich et~al.(2024)]{russinovich2024}
M.~Russinovich, A.~Salem, and R.~Eldan.
\newblock Great, now write an article about that: The crescendo multi-turn LLM jailbreak attack.
\newblock arXiv preprint, 2024.

\bibitem[SafeDream(2025)]{safedream2025}
SafeDream: Proactive safety for multi-turn conversations.
\newblock arXiv preprint, 2025.

\bibitem[Searle(1969)]{searle1969}
J.~R. Searle.
\newblock \emph{Speech Acts: An Essay in the Philosophy of Language}.
\newblock Cambridge University Press, 1969.

\bibitem[Stolcke et~al.(2000)]{stolcke2000}
A.~Stolcke, K.~Ries, N.~Coccaro, E.~Shriberg, R.~Bates, D.~Jurafsky, P.~Taylor, R.~Martin, C.~Van~Ess-Dykema, and M.~Meteer.
\newblock Dialogue act modeling for automatic tagging and recognition of conversational speech.
\newblock \emph{Computational Linguistics}, 26(3):339--373, 2000.

\bibitem[Su et~al.(2026)]{su2026}
Y.~Su et~al.
\newblock DeepContext: Multi-turn jailbreak detection via temporal semantic drift.
\newblock arXiv:2602.16935, 2026.

\bibitem[Tan et~al.(2016)]{tan2016}
C.~Tan, V.~Niculae, C.~Danescu-Niculescu-Mizil, and L.~Lee.
\newblock Winning arguments: Interaction dynamics and persuasion strategies in good-faith online discussions.
\newblock In \emph{WWW}, 2016.

\bibitem[Wei et~al.(2024)]{wei2024}
A.~Wei, N.~Haghtalab, and J.~Steinhardt.
\newblock Jailbroken: How does LLM safety training fail?
\newblock In \emph{NeurIPS}, 2024.

\bibitem[Zeng et~al.(2024a)]{zeng2024pap}
Y.~Zeng, H.~Lin, J.~Zhang, D.~Yang, R.~Jia, and W.~Shi.
\newblock How johnny can persuade LLMs to jailbreak them: Rethinking persuasion to challenge AI safety by humanizing LLMs.
\newblock In \emph{ACL}, 2024.

\bibitem[Zeng et~al.(2024b)]{zeng2024shieldgemma}
W.~Zeng et~al.
\newblock ShieldGemma: Generative AI content moderation based on Gemma.
\newblock arXiv preprint, 2024.

\bibitem[Zhao et~al.(2024)]{zhao2024wildchat}
W.~Zhao et~al.
\newblock WildChat: 1M ChatGPT interaction logs in the wild.
\newblock In \emph{ICLR}, 2024.

\end{thebibliography}

\end{document}
